%%%%%%%%%%%%%%%%%%%%%%%%%%%%%%%%%%%%%%%%%%%%%%%%%%%%%%%%%%%%
\section{Background}
\label{sec:background}
%%%%%%%%%%%%%%%%%%%%%%%%%%%%%%%%%%%%%%%%%%%%%%%%%%%%%%%%%%%%

\subsection{Chemistry Foundational Models}

\molm{} is an encoder-only masked language model pre-trained on 1.1 billion molecules from PubChem and ZINC via a masked language modeling objective. Its downstream applications include molecular property prediction and multimodal battery electrolyte formulation performance prediction \citep{ross2022largescalechemicallanguagerepresentations}. \smi{} is a 289-million-parameter encoder-decoder foundation model pre-trained on 91 million curated SMILES from PubChem through a combined masked language modeling and token reconstruction objective, evaluated on molecular property classification and regression across eleven MoleculeNet benchmarks \citep{wu2018moleculenet}, SMILES reconstruction, and battery electrolyte property prediction \citep{soares2024smi}.

\para{Chemical Properties and SMILES String}

To probe learned structural representations, each atom is annotated with eight RDKit-derived properties. A key distinction exists between properties that are determined directly by SMILES tokenization and those that require contextual computation. Atom type and aromaticity fall into the first category, as element symbols are explicitly encoded as distinct tokens, and aromatic atoms are marked by lowercase letters in SMILES notation, meaning these properties are fully recoverable from the embedding lookup before any transformer computation takes place. By contrast, hybridization, degree, total valence, ring membership, and chiral tag cannot be read from individual tokens in isolation and require the model to integrate information across the molecular graph, making them genuine tests of what the encoder has learned beyond its tokenization. Detailed definitions of all properties are provided in Appendix~\ref{appendix:properties}.

% \paragraph{Chemical property labels.}

\section{Related Work}

\para{Interpretability of Chemical Language Models}

Prior interpretability analyses of chemical language models have operated at the level of output embeddings or input-output correlations, without examining internal layer-by-layer representations. \citet{soares2024smi} construct a dataset of six carbon-chain molecular families, apply t-SNE to visualize the final embedding space, and fit linear regressions to SMILES embedding triplets to test whether molecular composition is linearly represented in the latent space. \smi{} achieves $R^2 = 0.99$ on this task compared to $R^2 = 0.55$ for \molm{}, suggesting that \smi{} organizes molecular information more linearly. However, this analysis is confined to the output embedding layer and does not examine how representations evolve across encoder depth or interrogate the attention mechanism. \citet{sharma2025foundation} define SMARTS strings for over 550 chemical substructures, compute per-molecule functional-group presence vectors aggregated into formulation descriptors weighted by molar composition, and apply Spearman correlation to relate chemical functional groups to downstream prediction performance. This approach characterizes relationships between chemical inputs and model outputs but does not inspect any internal model representations. In both cases, the analyses remain at the surface of the model and do not address which layers encode which properties or whether the encoded features are causally used in computation.

\para{Mechanistic Interpretability Tools from NLP}

The mechanistic interpretability tools we employ have well-established foundations in NLP research. Linear classifier probes were introduced by \citet{alain2016understanding}, who proposed fitting independent linear classifiers to the hidden states at each layer of a neural network and demonstrated that probe accuracy increases monotonically with depth, reflecting how linearly separable features become as the network transforms its input. \citet{tenney2019bert} applied this framework to BERT using edge-probing tasks derived from the classical NLP pipeline, showing that syntactic information is encoded in earlier layers while semantic information emerges in higher layers, and that the model implicitly recapitulates the traditional hierarchy of linguistic processing. \citet{clark2019what} complemented this representational analysis by directly examining BERT's attention weights, finding that individual attention heads specialize in specific syntactic relations and that attention maps capture substantial syntactic structure without explicit syntactic supervision. The linear attention mechanism used by both \molm{} and \smi{} is based on the FAVOR$+$ approach introduced by \citet{choromanski2021rethinking}, which approximates softmax attention via positive orthogonal random features to achieve linear time and space complexity. Because this mechanism does not form an explicit attention matrix, reconstructing effective attention weights for \smi{} requires extracting query and key tensors and applying the feature map directly, as detailed in Section~\ref{sec:methods}.

Our work brings mechanistic interpretability tools to chemical language models, establishing causal evidence for how structural properties are encoded. Rather than examining final-layer embeddings or input-output correlations, we track representations across all twelve encoder layers and separately distinguish properties recoverable from SMILES tokenization from those requiring genuine contextual computation. By combining linear probing with causal intervention, we move beyond demonstrating correlation to establishing that probed directions are actively maintained and used by the encoder. This allows us to draw a sharper distinction between \molm{} and \smi{}. Both models exhibit correlational signals for contextual chemical properties, but only \smi{} organizes these signals as stable, causally-used linear directions that propagate robustly through subsequent computation.

% \subsection{SMILES as a Molecular Language}
% \label{sec:smiles}

% The Simplified Molecular-Input Line Entry System (SMILES) encodes a molecule as a
% character string by performing a depth-first pre-order traversal of the molecular
% graph, generating symbols for atoms, bonds, branching decisions, and ring closures. For example, ethanol is written as \texttt{CCO} and benzene is written
% as \texttt{c1ccccc1}, where lowercase letters are used for aromatic atoms and digits indicate
% ring closures. SMILES is widely adopted because it is compact relative to
% graph representations and explicitly encodes meaningful substructures such as branches,
% cyclic systems, and chirality. However, it does not directly encode 3D geometry or
% interatomic distances, which are crucial for predicting quantum-mechanical properties.
% A central question in chemical language modeling is therefore whether models trained
% on SMILES alone can implicitly recover geometric information.

% %%%%%%%%%%%%%%%%%%%%%%%%%%%%%%%%%%%%%%%%%%%%%%%%%%%%%%%%%%%%
% \subsection{MoLFormer-XL}
% \label{sec:molformer}
% %%%%%%%%%%%%%%%%%%%%%%%%%%%%%%%%%%%%%%%%%%%%%%%%%%%%%%%%%%%%

% \subsubsection{Architecture}
% \label{sec:molformer-arch}

% MoLFormer-XL is an \emph{encoder-only} transformer with 12 layers,
% 12 attention heads per layer, and a hidden dimension of 768, giving approximately
% 47~million parameters in the encoder. It uses a vocabulary of 2{,}362 tokens
% (2{,}357 unique SMILES tokens plus 5 special tokens), with a maximum sequence length of
% 202 tokens (sufficient for 99.4\% of PubChem and ZINC molecules).

% MoLFormer-XL uses rotary linear attention with a generalized random feature
% map $\phi$:
% \begin{equation*}
% \small
% \mathrm{Attention}_m(Q,K,V) =
% \frac{\sum_{n=1}^{N} \langle \phi(R_m q_m),\, \phi(R_n k_n) \rangle v_n}
%      {\sum_{n=1}^{N} \langle \phi(R_m q_m),\, \phi(R_n k_n) \rangle}
% \end{equation*}
% where $R_m$ denotes the rotary position matrix at position $m$.

% After processing the sequence, the model produces a hidden representation $h_i$
% for each token. A single molecular embedding is obtained via mean pooling:
% \begin{equation*}
%   h_{\mathrm{mol}} = \frac{1}{N} \sum_{i=1}^{N} h_i.
% \end{equation*}
% This mean pooled vector is used for downstream prediction tasks.

% \paragraph{Implications for structure and geometry.}
% Although SMILES does not explicitly encode three-dimensional structure, rotary
% position encoding enables consistent relative-position representations, while
% attention captures interactions between chemically related tokens. Together,
% these mechanisms allow the model to implicitly encode aspects of spatial structure
% and approximate interatomic relationships from sequence alone. \xiaoyan{We should highlight what they didn't do. They just did a brief correlation analysis. We can add the causality. }

% \subsubsection{Training Data and Protocol}
% \label{sec:molformer-training}

% MoLFormer-XL is pretrained on \textbf{1.1 billion molecules}: 111~million from
% PubChem and approximately 1~billion from ZINC, all canonicalised via RDKit.
% Pretraining is performed using masked language modelling (MLM) in the BERT
% tradition, with the objective of reconstructing the original SMILES string from
% partially observed inputs. Specifically, 15\% of tokens in each sequence are selected
% at random; of these, 80\% are replaced with the special \texttt{[MASK]} token,
% 10\% are replaced with a random token, and 10\% are left unchanged. The model is
% then trained to predict the original token at each selected position using the
% surrounding context.

% This training objective forces the model to learn the syntactic and semantic
% structure of SMILES strings, including local bonding patterns, branching, and
% long-range dependencies such as ring closures. Importantly, no explicit supervision
% on molecular properties or geometric structure is provided during pretraining;
% all learned representations arise from this token-level reconstruction task. 

% %%%%%%%%%%%%%%%%%%%%%%%%%%%%%%%%%%%%%%%%%%%%%%%%%%%%%%%%%%%%
% \subsection{SMI-TED$_{289\mathrm{M}}$}
% \label{sec:smited}
% %%%%%%%%%%%%%%%%%%%%%%%%%%%%%%%%%%%%%%%%%%%%%%%%%%%%%%%%%%%%

% \subsubsection{Architecture}
% \label{sec:smited-arch}

% SMI-TED$_{289\mathrm{M}}$ is a chemical foundation
% model family with an \textit{encoder-decoder} architecture. The encoder is a
% 12-layer, 12-head, 768-dimensional bidirectional transformer (47M parameters). SMI-TED$_{289\mathrm{M}}$
% also uses rotary linear attention with a generalized random feature map $\phi$.

% The decoder is a bottleneck
% autoencoder (242M parameters) that compresses the full token embedding matrix into a
% fixed-size molecular embedding and enables reconstruction of the original SMILES.

% Formally, given token embeddings $x \in \mathbb{R}^{D \times L}$ (where $D = 202$
% is the maximum token count and $L = 768$ is the embedding dimension), the encoder
% produces a molecular latent vector $z \in \mathbb{R}^L$ via:
% \begin{equation*}
%   z = \bigl(\mathrm{LayerNorm}(\mathrm{GELU}(x W_1 + b_1))\bigr) W_2,
%   \label{eq:smited-enc}
% \end{equation*}
% with $W_1 \in \mathbb{R}^{D \times L}$, $b_1 \in \mathbb{R}^L$,
% $W_2 \in \mathbb{R}^{L \times L}$. The latent vector $z$ can be decoded back to
% approximate token embeddings $\hat{x}$ via a symmetric linear-GELU-LayerNorm
% projection, followed by a language model head that reconstructs individual tokens.

% This \emph{invertible} latent space is the key structural difference from
% MoLFormer-XL’s mean-pooled representation. The encoder--decoder architecture imposes
% a reconstruction constraint: the latent vector $z$ must retain sufficient
% information to recover the original SMILES sequence, rather than allowing
% information to be discarded during pooling. This bottleneck encourages a compact
% yet information-preserving representation that captures global molecular
% structure. As a result, SMI-TED$_{289\mathrm{M}}$ learns more structured and compositional latent
% embeddings, which are better suited for probing and interpretability analysis.

% \subsubsection{Training Data and Protocol}
% \label{sec:smited-training}

% SMI-TED$_{289\mathrm{M}}$ is pretrained on \textbf{91 million} carefully curated molecules from
% PubChem. 

% \textbf{Phase 1} (95/5 split): The token encoder is pretrained on 95\% of the data
% using MLM (same 15\%:80\%/10\%/10\% masking recipe as MoLFormer-XL). The remaining 5\%
% is used for the encoder-decoder reconstruction loss. This separation is necessary
% because unstable token embeddings in early epochs would otherwise destabilise decoder
% training.

% \textbf{Phase 2} (100\% data): Once the encoder has converged, both the MLM loss
% (encoder) and the reconstruction loss (encoder-decoder) are jointly optimised on
% the full dataset, improving decoder performance.